% =====================================================================
% Spacecraft Attitude Dynamics and Control — Session Summary
%
% USAGE
%   Standalone:  compile this .tex directly (it's a complete document).
%   Embedded:    strip the preamble (everything before \begin{document})
%                and % =====================================================================
% Spacecraft Attitude Dynamics and Control — Session Summary
%
% USAGE
%   Standalone:  compile this .tex directly (it's a complete document).
%   Embedded:    strip the preamble (everything before \begin{document})
%                and % =====================================================================
% Spacecraft Attitude Dynamics and Control — Session Summary
%
% USAGE
%   Standalone:  compile this .tex directly (it's a complete document).
%   Embedded:    strip the preamble (everything before \begin{document})
%                and % =====================================================================
% Spacecraft Attitude Dynamics and Control — Session Summary
%
% USAGE
%   Standalone:  compile this .tex directly (it's a complete document).
%   Embedded:    strip the preamble (everything before \begin{document})
%                and \input{session_summary.tex} into your existing project.
%
% BIBLIOGRAPHY
%   The single \cite{markley2014} in this document needs the following
%   entry in your .bib file:
%
%   @book{markley2014,
%     author    = {Markley, F. Landis and Crassidis, John L.},
%     title     = {Fundamentals of Spacecraft Attitude Determination and Control},
%     publisher = {Springer},
%     year      = {2014},
%     series    = {Space Technology Library},
%     volume    = {33},
%     address   = {New York, NY},
%     doi       = {10.1007/978-1-4939-0802-8}
%   }
%
%   If you don't use BibTeX, find-and-replace \cite{markley2014} with
%   "Markley \& Crassidis (2014)".
% =====================================================================

\documentclass[11pt]{article}

\usepackage[margin=1in]{geometry}
\usepackage{amsmath,amssymb,bm}
\usepackage{listings}
\usepackage{xcolor}
\usepackage{booktabs}
\usepackage{enumitem}
\usepackage{hyperref}
\usepackage{tcolorbox}

\hypersetup{colorlinks=true, linkcolor=blue!60!black, urlcolor=blue!60!black}

\lstdefinestyle{pycode}{
    language=Python,
    basicstyle=\ttfamily\footnotesize,
    keywordstyle=\color{blue!70!black}\bfseries,
    stringstyle=\color{teal!70!black},
    commentstyle=\color{gray}\itshape,
    showstringspaces=false,
    breaklines=true,
    frame=single,
    rulecolor=\color{black!20},
    tabsize=4,
    columns=flexible,
    aboveskip=8pt,
    belowskip=8pt,
}
\lstset{style=pycode}

\title{Spacecraft Attitude Dynamics and Control Simulation\\[3pt]
       \large Session Build-Up Summary}
\author{}
\date{\today}


\begin{document}
\maketitle


% ---------------------------------------------------------------------
\section{Project Goal}
% ---------------------------------------------------------------------

Build, incrementally, a self-contained Python simulation of spacecraft
attitude dynamics and control following Markley and
Crassidis~\cite{markley2014}. Each step in the session added a piece of
physics or refactored the code for clarity. The final product is two
deliverables --- a working simulation in Python and a companion
LaTeX summary of the underlying physics --- plus an automated smoke test.


% ---------------------------------------------------------------------
\section{System Overview}
% ---------------------------------------------------------------------

The final simulation models a rigid spacecraft with the following
components:

\begin{itemize}\setlength{\itemsep}{2pt}
  \item Quaternion-feedback PD attitude controller (scalar-last convention).
  \item Four reaction wheels in a NASA-standard pyramid configuration.
  \item Magnetic torquers performing wheel momentum unloading via a
        cross-product (``\(B\)-dot--style'') control law.
  \item A propellant-slosh disturbance modelled as a forced second-order
        ODE on a 3-D torque state.
  \item A 3-2-1 (Z-Y-X) Euler-angle interface for human-readable initial
        conditions and post-processing.
\end{itemize}

\noindent The integrator is \texttt{scipy.integrate.solve\_ivp} (adaptive
RK45 by default), and the closed-loop is treated as continuous-time --- the
PD, magnetorquer, and wheel-allocation laws are evaluated at every solver
step rather than held with zero-order hold.


% ---------------------------------------------------------------------
\section{Modelling Components}
% ---------------------------------------------------------------------

\subsection{Quaternion Kinematics (Scalar-Last)}

Using \(\bm{q}=[\bm{q}_v;\,q_4]\), the body-to-inertial quaternion satisfies
\begin{equation}
\dot{\bm{q}} = \tfrac{1}{2}\,\bm{\Omega}(\bm{\omega})\,\bm{q}
\;=\;
\begin{bmatrix}
 \tfrac{1}{2}\bigl(q_4\,\bm{\omega} - \bm{\omega}\times\bm{q}_v\bigr) \\
-\tfrac{1}{2}\,\bm{\omega}^{\!\top}\bm{q}_v
\end{bmatrix},
\tag{M\&C Eq.~2.88}
\end{equation}
and the quaternion product follows
\(\bm{A}(\bm{q}\otimes\bm{p}) = \bm{A}(\bm{q})\bm{A}(\bm{p})\)
(M\&C~Eq.~2.82b).

\subsection{Rigid-Body Dynamics with Reaction Wheels}

Euler's equation with internal angular momentum
\(\bm{h}_w^{\,b} = \bm{W}\bm{h}_w\) is
\begin{equation}
\bm{J}\,\dot{\bm{\omega}}
\;=\;
\bm{u}_{\text{body}}
\;-\;
\bm{\omega}\times(\bm{J}\bm{\omega} + \bm{h}_w^{\,b}),
\tag{M\&C Sec.~7.2}
\end{equation}
where the inertia tensor uses the sign convention of M\&C~Eq.~3.13.
Wheel torque commands integrate the wheel-momentum law
\(\dot{\bm{h}}_w = \bm{\tau}_w\).

\subsection{Reaction-Wheel Allocation}

A four-wheel NASA pyramid with cant angle \(\beta = \arccos(1/\sqrt{3})
\approx 54.74^\circ\) gives equal projection of each wheel onto the body
axes. A desired body torque is allocated by the minimum-norm pseudoinverse:
\begin{equation}
\bm{\tau}_w \;=\; -\bm{W}^{+}\,\bm{u}_{\text{cmd}},
\qquad
\bm{W}^{+} = \bm{W}^{\!\top}(\bm{W}\bm{W}^{\!\top})^{-1}.
\end{equation}

\subsection{Quaternion-Feedback PD Controller}

With error quaternion \(\delta\bm{q} = \bm{q}_{\text{des}}^{-1}\otimes\bm{q}\),
the controller (M\&C~Sec.~7.4) is
\begin{equation}
\bm{u}_{\text{cmd}} \;=\;
\operatorname{sat}_{u_{\max}}\!\Bigl(
 -K_p\,\operatorname{sign}(\delta q_4)\,\delta\bm{q}_v
 \;-\; K_d\,\bm{\omega}
\Bigr).
\end{equation}
The \(\operatorname{sign}(\delta q_4)\) factor selects the shortest rotation
and prevents the well-known \emph{unwinding} phenomenon. Each component is
clipped to \([-u_{\max},+u_{\max}]\).

\subsection{Magnetic Momentum Unloading}

The cross-product (``\(B\)-dot--style'') law from M\&C~Sec.~7.5
\begin{equation}
\bm{m} \;=\; \operatorname{sat}_{m_{\max}}\!\biggl(
\frac{K_{\text{dump}}}{\|\bm{B}\|^{2}}\,
\bm{h}_w^{\,b}\times\bm{B}\biggr),
\quad
\bm{\tau}_{\text{mag}} \;=\; \bm{m}\times\bm{B},
\end{equation}
applies a damping torque to the component of \(\bm{h}_w^{\,b}\)
perpendicular to \(\bm{B}\). Orbital variation of \(\bm{B}\) sweeps the
unloadable subspace through all three body axes over one orbit. The toy
geomagnetic field is an inclined rotating dipole of magnitude
\(\sim 30~\mu\mathrm{T}\) at 90-min LEO period.

\subsection{Propellant Slosh Model}

Slosh torque \(\bm{T}_s \in \mathbb{R}^{3}\) is governed by
\begin{equation}
\ddot{\bm{T}}_s \;=\;
 -A_s(t)\,\bm{\omega}
 \;-\; B_s(t)\,\dot{\bm{\omega}}
 \;-\; C_s(t)\,\dot{\bm{T}}_s
 \;-\; K_s(t)\,\bm{T}_s,
\label{eq:slosh}
\end{equation}
with each scalar coefficient time-varying as
\begin{equation}
f(t) \;=\; f_{\text{nom}} + f_{\text{amp}}\sin(\omega_f\,t),
\qquad f \in \{A_s,B_s,C_s,K_s\}.
\end{equation}
The constraint \(f_{\text{nom}} > |f_{\text{amp}}|\) keeps \(K_s, C_s\)
positive (negative stiffness or damping would be unphysical). \(\bm{T}_s\)
enters Euler's equation as an additional external torque on the body.

\textbf{No algebraic loop.} \(\dot{\bm{\omega}}\) depends on \(\bm{T}_s\) (a
\emph{state}, not on \(\ddot{\bm{T}}_s\)), and \(\ddot{\bm{T}}_s\) consumes
the just-computed \(\dot{\bm{\omega}}\). The RHS is fully explicit.

\subsection{Coupled Equations of Motion}

The complete state is
\(\bm{x} = [\bm{q};\bm{\omega};\bm{h}_w;\bm{T}_s;\dot{\bm{T}}_s]
        \in \mathbb{R}^{13+n_w}\),
governed by
\begin{align}
\dot{\bm{q}}        &= \tfrac{1}{2}\,\bm{\Omega}(\bm{\omega})\,\bm{q}, \\
\bm{J}\,\dot{\bm{\omega}} &= \bm{u}_{\text{body}}
                            - \bm{\omega}\times(\bm{J}\bm{\omega}+\bm{W}\bm{h}_w), \\
\dot{\bm{h}}_w      &= \bm{\tau}_w, \\
\ddot{\bm{T}}_s     &= -A_s\bm{\omega} - B_s\dot{\bm{\omega}}
                       - C_s\dot{\bm{T}}_s - K_s\bm{T}_s,
\end{align}
with
\(\bm{u}_{\text{body}} = -\bm{W}\bm{\tau}_w + \bm{m}\times\bm{B} + \bm{T}_s\).


% ---------------------------------------------------------------------
\section{Code Architecture Evolution}
% ---------------------------------------------------------------------

The simulation was built incrementally; each step is reproducible from
the conversation log.

\begin{enumerate}\setlength{\itemsep}{4pt}
  \item \textbf{Baseline rigid body + PD controller.}
        A single-file implementation of Euler's equation with a quaternion
        PD law and a hand-written fourth-order Runge--Kutta integrator.
  \item \textbf{Euler-angle initial conditions.}
        Helper \texttt{euler321\_to\_quat()} converts a roll-pitch-yaw
        triple to a scalar-last quaternion (M\&C App.~B.96).
  \item \textbf{Full 3$\times$3 inertia tensor.}
        Replaced the diagonal default with off-diagonal products of inertia
        and added symmetry / positive-definiteness assertions.
  \item \textbf{Reaction-wheel dynamics.}
        Added the four-wheel NASA pyramid configuration matrix \(\bm{W}\),
        pseudoinverse allocation, and the wheel-momentum gyroscopic term in
        Euler's equation.
  \item \textbf{Magnetic momentum unloading.}
        Added a toy LEO B-field, the cross-product unloading law, and per-
        axis dipole saturation. Also introduced PD torque saturation
        \(u_{\max}\).
  \item \textbf{Propellant slosh disturbance.}
        Added the second-order slosh ODE driven by \((\bm{\omega},
        \dot{\bm{\omega}})\), with time-varying scalar coefficients
        (DC~+~sine).
  \item \textbf{Variable rename.}
        Renamed \(\Gamma_s \to T_s\) throughout the code and plot labels
        for consistency with the user's preferred notation.
  \item \textbf{Toggles.}
        Added independent on/off switches for slosh, magnetorquer, and
        reaction wheels. With wheels disabled the PD output drives the body
        directly as an ideal torque actuator.
  \item \textbf{Migration to \texttt{solve\_ivp}.}
        Replaced the hand-rolled RK4 loop with
        \texttt{scipy.integrate.solve\_ivp} (adaptive RK45). Quaternion
        re-normalisation moved from per-step to post-processing.
  \item \textbf{Merge and clean-up.}
        Consolidated the two intermediate files into a single
        \texttt{rigid\_body\_pd.py}. Removed all RK4-specific code paths.
  \item \textbf{Dataclass + presets refactor (v2).}
        Centralised every spacecraft parameter, gain, initial condition,
        time setting, and toggle into a single \texttt{SimConfig} dataclass
        at the top of the file. Added two named preset factories,
        \texttt{mc\_example\_71()} and \texttt{mc\_example\_72()}, matching
        the worked examples in M\&C Sec.~7.
\end{enumerate}


% ---------------------------------------------------------------------
\section{Final Architecture (v2)}
% ---------------------------------------------------------------------

\subsection{The \texttt{SimConfig} Dataclass}

All scenario knobs live in one place:

\begin{lstlisting}
@dataclass
class SimConfig:
    """All satellite parameters, gains, ICs, time, and toggles."""
    # Inertia (M&C Eq. 3.13), kg.m^2
    Ixx: float = 6400.0;  Iyy: float = 4730.0;  Izz: float = 8160.0
    Ixy: float = 76.4;    Ixz: float = 25.6;    Iyz: float = 40.0
    # Wheel geometry
    wheel_tilt_deg: float = 54.7356      # 4-wheel NASA pyramid
    # PD gains and saturation
    Kp: float = 10.0;   Kd: float = 150.0;   u_max: float = 1000.0
    # Magnetorquer
    K_dump: float = 1.0;  m_max: float = 500.0
    # Slosh:  f(t) = nom + amp*sin(omega*t)
    slosh_params: dict = field(default_factory=lambda: {
        'A': {'nom': 0.005, 'amp': 0.0, 'omega': 0.05},
        'B': {'nom': 0.020, 'amp': 0.0, 'omega': 0.05},
        'C': {'nom': 0.023, 'amp': 0.0, 'omega': 0.05},
        'K': {'nom': 0.008, 'amp': 0.0, 'omega': 0.05},
    })
    # Initial conditions (scalar-last)
    q0:    np.ndarray = field(default_factory=lambda:
        (math.sqrt(2)/2) * np.array([1.0, 0.0, 0.0, 1.0]))
    w0:    np.ndarray = field(default_factory=lambda:
        np.array([0.01, 0.01, 0.01]))
    q_des: np.ndarray = field(default_factory=lambda:
        np.array([0.0, 0.0, 0.0, 1.0]))
    Ts0:   np.ndarray = field(default_factory=lambda: np.zeros(3))
    Tsd0:  np.ndarray = field(default_factory=lambda: np.zeros(3))
    # Time
    t_end: float = 20*60.0;  dt_eval: float = 0.05
    # Toggles
    enable_slosh:        bool = False
    enable_magnetorquer: bool = False
    enable_wheels:       bool = True
    # Solver
    method: str = 'RK45';  rtol: float = 1e-8
    atol: float = 1e-10;   max_step: float = 0.5
\end{lstlisting}

\subsection{Named Presets}

\begin{lstlisting}
def mc_example_71() -> SimConfig:
    """Large spacecraft, diagonal J, fast initial tumble."""
    cfg = SimConfig()
    cfg.Ixx, cfg.Iyy, cfg.Izz = 10000.0, 9000.0, 12000.0
    cfg.Ixy = cfg.Ixz = cfg.Iyz = 0.0
    cfg.Kp, cfg.Kd = 50.0, 500.0
    cfg.u_max = 100.0
    q = np.array([0.6853, 0.6953, 0.1531, 0.1531])
    cfg.q0 = q / np.linalg.norm(q)
    cfg.w0 = np.array([0.53, 0.53, 0.053])
    cfg.t_end = 300.0
    return cfg

def mc_example_72() -> SimConfig:
    """Smaller spacecraft with products of inertia; gentle IC."""
    return SimConfig()    # SimConfig() defaults already match Example 7.2
\end{lstlisting}

\subsection{Usage Pattern}

\begin{lstlisting}
cfg = mc_example_72()
cfg.enable_slosh = True            # in-place tweak
cfg.Kp, cfg.Kd = 15.0, 200.0
results = simulate(cfg)            # one entry point
plot_results(*results)
plot_euler_and_rates(results[0], results[1], results[3])
plt.show()
\end{lstlisting}


% ---------------------------------------------------------------------
\section{Validation: Automated Smoke Test}
% ---------------------------------------------------------------------

A separate file \texttt{test\_rigid\_body\_pd\_v2.py} exercises the
simulation under every toggle combination using the non-interactive
\texttt{Agg} matplotlib backend. All eight test cases pass.

\begin{center}
\begin{tabular}{lc}
\toprule
\textbf{Test} & \textbf{Result} \\
\midrule
Euler $\leftrightarrow$ quaternion roundtrip (4 angle triples) & Pass \\
Example 7.1 preset, wheels only                                 & Pass \\
Example 7.2 preset, wheels only                                 & Pass \\
Example 7.2 full stack (slosh + mag + wheels)                   & Pass \\
Example 7.2 with wheels OFF (ideal actuator)                    & Pass \\
Example 7.2 with all subsystems OFF                             & Pass \\
Example 7.2 with overridden gains (in-place mutation)           & Pass \\
Plot functions (both figures)                                   & Pass \\
\bottomrule
\end{tabular}
\end{center}

\noindent Per-trajectory invariants checked at every output sample:
\begin{itemize}\setlength{\itemsep}{1pt}
  \item Output tuple length = 11; all array shapes match the expected
        \((N,4)\), \((N,3)\), \((N,n_w)\).
  \item No NaNs in any state trajectory.
  \item Unit-norm quaternion preserved to machine precision
        (\(\max\bigl|\,\|\bm{q}\|-1\,\bigr| = 2.22\times10^{-16}\)).
  \item \(|\bm{\omega}|\) monotonically non-increasing in every settling case.
  \item When a toggle is OFF, the corresponding telemetry stays at
        exactly zero (\(\bm{h}_w\), \(\bm{\tau}_w\), \(\bm{m}\),
        \(\bm{\tau}_{\text{mag}}\), \(\bm{T}_s\)).
\end{itemize}


% ---------------------------------------------------------------------
\section{Companion LaTeX Physics Summary}
% ---------------------------------------------------------------------

A separate file \texttt{attitude\_dynamics\_model.tex} was produced for
direct import into the project's Overleaf write-up. It contains a
self-contained section with every equation above plus extended
discussion, fully cross-referenced via \texttt{\textbackslash label}.
The \texttt{\textbackslash cite\{markley2014\}} entry should resolve to
Markley and Crassidis using the BibTeX stanza included at the top of that
file.


% ---------------------------------------------------------------------
\section{Proposed Extension: NARX Slosh Surrogate}
% ---------------------------------------------------------------------

A design (no code) was discussed for replacing the analytical slosh ODE
\eqref{eq:slosh} with a learned Nonlinear Autoregressive Network with
eXogenous Inputs (NARX). The map to learn is
\begin{equation}
\hat{T}_s(k) \;=\; f_\theta\bigl(
  T_s(k{-}1),\dots,T_s(k{-}n_y),\;
  \omega(k),\dots,\omega(k{-}n_u),\;
  \dot\omega(k),\dots,\dot\omega(k{-}n_u)
\bigr),
\end{equation}
optionally extended to a \emph{parametric} surrogate by also feeding the
four slosh coefficients \((A_s, B_s, C_s, K_s)\). The plan, in order:

\begin{enumerate}\setlength{\itemsep}{2pt}
  \item Add a helper \texttt{reconstruct\_w\_dot()} that recomputes
        \(\dot{\bm{\omega}}\) exactly from the saved \texttt{simulate()}
        outputs (no finite differences).
  \item Implement \texttt{sample\_random\_config(rng)} that draws
        \texttt{SimConfig} instances with randomised initial conditions,
        gains, inertia perturbations, and (optionally) slosh coefficients.
  \item Generate 100--500 trajectories, each $\sim$600\,s, downsampled to
        \(\Delta t \approx 0.5\,\mathrm{s}\) and stacked across all three
        body axes (the slosh dynamics are axis-decoupled, so one scalar
        NARX trained on stacked per-axis data is the right architecture).
  \item Build sliding-window features with lag orders
        \(n_y \approx n_u \approx 10\)--\(20\) samples.
  \item Train a small MLP (e.g. 64-64 \texttt{tanh}) with
        \texttt{sklearn.neural\_network.MLPRegressor} as a baseline, then
        promote to PyTorch if needed.
  \item Validate in two modes:
        \emph{one-step prediction} (teacher-forced lags) for training-loss
        sanity, and \emph{free-running rollout} (model feeds its own past
        predictions) for honest dynamic accuracy.
  \item Split by \emph{whole trajectory} (70/15/15) to avoid leakage from
        adjacent correlated samples within one simulation.
  \item Optional integration: add \texttt{cfg.use\_narx\_slosh = True}
        and \texttt{cfg.narx\_model = trained\_model} to the dataclass,
        wire it into \texttt{dynamics\_rhs}, and compare closed-loop
        traces against the analytical slosh ground truth on held-out
        initial conditions.
\end{enumerate}


% ---------------------------------------------------------------------
\section{Deliverables}
% ---------------------------------------------------------------------

\begin{center}
\begin{tabular}{ll}
\toprule
\textbf{File} & \textbf{Purpose} \\
\midrule
\texttt{rigid\_body\_pd.py}            & Merged simulation, \texttt{solve\_ivp}-based \\
\texttt{rigid\_body\_pd\_v2.py}        & Dataclass + presets refactor (recommended) \\
\texttt{test\_rigid\_body\_pd\_v2.py}  & Smoke test (8 cases, all pass) \\
\texttt{attitude\_dynamics\_model.tex} & Physics section for Overleaf \\
\texttt{session\_summary.tex}          & This document \\
\bottomrule
\end{tabular}
\end{center}


% ---------------------------------------------------------------------
\section{Conventions Used Throughout}
% ---------------------------------------------------------------------

\begin{itemize}\setlength{\itemsep}{1pt}
  \item \textbf{Quaternion:} scalar-last, \(\bm{q}=[q_1, q_2, q_3, q_4]\)
        with \(q_4\) the scalar part.
  \item \textbf{Euler sequence:} 3-2-1 (Z-Y-X, yaw-pitch-roll) for human
        I/O; quaternions everywhere internally.
  \item \textbf{Inertia sign:} off-diagonal entries are stored as positive
        products of inertia; the negative signs live in the \(\bm{J}\)
        builder (M\&C Eq.~3.13).
  \item \textbf{Frames:} body frame throughout; \(\bm{B}\) and
        \(\bm{h}_w^{\,b}\) are body-frame projections.
  \item \textbf{Slosh torque symbol:} \(\bm{T}_s\) (code: \texttt{Ts}),
        with rate \(\dot{\bm{T}}_s\) (\texttt{Ts\_dot}). Renamed from the
        original \(\Gamma_s\).
\end{itemize}


% ---------------------------------------------------------------------
\begin{thebibliography}{1}
% ---------------------------------------------------------------------

\bibitem{markley2014}
F.~L.~Markley and J.~L.~Crassidis,
\emph{Fundamentals of Spacecraft Attitude Determination and Control},
Space Technology Library, Vol.~33,
Springer, New York, NY, 2014.
\href{https://doi.org/10.1007/978-1-4939-0802-8}{doi:10.1007/978-1-4939-0802-8}.

\end{thebibliography}

\end{document}
 into your existing project.
%
% BIBLIOGRAPHY
%   The single \cite{markley2014} in this document needs the following
%   entry in your .bib file:
%
%   @book{markley2014,
%     author    = {Markley, F. Landis and Crassidis, John L.},
%     title     = {Fundamentals of Spacecraft Attitude Determination and Control},
%     publisher = {Springer},
%     year      = {2014},
%     series    = {Space Technology Library},
%     volume    = {33},
%     address   = {New York, NY},
%     doi       = {10.1007/978-1-4939-0802-8}
%   }
%
%   If you don't use BibTeX, find-and-replace \cite{markley2014} with
%   "Markley \& Crassidis (2014)".
% =====================================================================

\documentclass[11pt]{article}

\usepackage[margin=1in]{geometry}
\usepackage{amsmath,amssymb,bm}
\usepackage{listings}
\usepackage{xcolor}
\usepackage{booktabs}
\usepackage{enumitem}
\usepackage{hyperref}
\usepackage{tcolorbox}

\hypersetup{colorlinks=true, linkcolor=blue!60!black, urlcolor=blue!60!black}

\lstdefinestyle{pycode}{
    language=Python,
    basicstyle=\ttfamily\footnotesize,
    keywordstyle=\color{blue!70!black}\bfseries,
    stringstyle=\color{teal!70!black},
    commentstyle=\color{gray}\itshape,
    showstringspaces=false,
    breaklines=true,
    frame=single,
    rulecolor=\color{black!20},
    tabsize=4,
    columns=flexible,
    aboveskip=8pt,
    belowskip=8pt,
}
\lstset{style=pycode}

\title{Spacecraft Attitude Dynamics and Control Simulation\\[3pt]
       \large Session Build-Up Summary}
\author{}
\date{\today}


\begin{document}
\maketitle


% ---------------------------------------------------------------------
\section{Project Goal}
% ---------------------------------------------------------------------

Build, incrementally, a self-contained Python simulation of spacecraft
attitude dynamics and control following Markley and
Crassidis~\cite{markley2014}. Each step in the session added a piece of
physics or refactored the code for clarity. The final product is two
deliverables --- a working simulation in Python and a companion
LaTeX summary of the underlying physics --- plus an automated smoke test.


% ---------------------------------------------------------------------
\section{System Overview}
% ---------------------------------------------------------------------

The final simulation models a rigid spacecraft with the following
components:

\begin{itemize}\setlength{\itemsep}{2pt}
  \item Quaternion-feedback PD attitude controller (scalar-last convention).
  \item Four reaction wheels in a NASA-standard pyramid configuration.
  \item Magnetic torquers performing wheel momentum unloading via a
        cross-product (``\(B\)-dot--style'') control law.
  \item A propellant-slosh disturbance modelled as a forced second-order
        ODE on a 3-D torque state.
  \item A 3-2-1 (Z-Y-X) Euler-angle interface for human-readable initial
        conditions and post-processing.
\end{itemize}

\noindent The integrator is \texttt{scipy.integrate.solve\_ivp} (adaptive
RK45 by default), and the closed-loop is treated as continuous-time --- the
PD, magnetorquer, and wheel-allocation laws are evaluated at every solver
step rather than held with zero-order hold.


% ---------------------------------------------------------------------
\section{Modelling Components}
% ---------------------------------------------------------------------

\subsection{Quaternion Kinematics (Scalar-Last)}

Using \(\bm{q}=[\bm{q}_v;\,q_4]\), the body-to-inertial quaternion satisfies
\begin{equation}
\dot{\bm{q}} = \tfrac{1}{2}\,\bm{\Omega}(\bm{\omega})\,\bm{q}
\;=\;
\begin{bmatrix}
 \tfrac{1}{2}\bigl(q_4\,\bm{\omega} - \bm{\omega}\times\bm{q}_v\bigr) \\
-\tfrac{1}{2}\,\bm{\omega}^{\!\top}\bm{q}_v
\end{bmatrix},
\tag{M\&C Eq.~2.88}
\end{equation}
and the quaternion product follows
\(\bm{A}(\bm{q}\otimes\bm{p}) = \bm{A}(\bm{q})\bm{A}(\bm{p})\)
(M\&C~Eq.~2.82b).

\subsection{Rigid-Body Dynamics with Reaction Wheels}

Euler's equation with internal angular momentum
\(\bm{h}_w^{\,b} = \bm{W}\bm{h}_w\) is
\begin{equation}
\bm{J}\,\dot{\bm{\omega}}
\;=\;
\bm{u}_{\text{body}}
\;-\;
\bm{\omega}\times(\bm{J}\bm{\omega} + \bm{h}_w^{\,b}),
\tag{M\&C Sec.~7.2}
\end{equation}
where the inertia tensor uses the sign convention of M\&C~Eq.~3.13.
Wheel torque commands integrate the wheel-momentum law
\(\dot{\bm{h}}_w = \bm{\tau}_w\).

\subsection{Reaction-Wheel Allocation}

A four-wheel NASA pyramid with cant angle \(\beta = \arccos(1/\sqrt{3})
\approx 54.74^\circ\) gives equal projection of each wheel onto the body
axes. A desired body torque is allocated by the minimum-norm pseudoinverse:
\begin{equation}
\bm{\tau}_w \;=\; -\bm{W}^{+}\,\bm{u}_{\text{cmd}},
\qquad
\bm{W}^{+} = \bm{W}^{\!\top}(\bm{W}\bm{W}^{\!\top})^{-1}.
\end{equation}

\subsection{Quaternion-Feedback PD Controller}

With error quaternion \(\delta\bm{q} = \bm{q}_{\text{des}}^{-1}\otimes\bm{q}\),
the controller (M\&C~Sec.~7.4) is
\begin{equation}
\bm{u}_{\text{cmd}} \;=\;
\operatorname{sat}_{u_{\max}}\!\Bigl(
 -K_p\,\operatorname{sign}(\delta q_4)\,\delta\bm{q}_v
 \;-\; K_d\,\bm{\omega}
\Bigr).
\end{equation}
The \(\operatorname{sign}(\delta q_4)\) factor selects the shortest rotation
and prevents the well-known \emph{unwinding} phenomenon. Each component is
clipped to \([-u_{\max},+u_{\max}]\).

\subsection{Magnetic Momentum Unloading}

The cross-product (``\(B\)-dot--style'') law from M\&C~Sec.~7.5
\begin{equation}
\bm{m} \;=\; \operatorname{sat}_{m_{\max}}\!\biggl(
\frac{K_{\text{dump}}}{\|\bm{B}\|^{2}}\,
\bm{h}_w^{\,b}\times\bm{B}\biggr),
\quad
\bm{\tau}_{\text{mag}} \;=\; \bm{m}\times\bm{B},
\end{equation}
applies a damping torque to the component of \(\bm{h}_w^{\,b}\)
perpendicular to \(\bm{B}\). Orbital variation of \(\bm{B}\) sweeps the
unloadable subspace through all three body axes over one orbit. The toy
geomagnetic field is an inclined rotating dipole of magnitude
\(\sim 30~\mu\mathrm{T}\) at 90-min LEO period.

\subsection{Propellant Slosh Model}

Slosh torque \(\bm{T}_s \in \mathbb{R}^{3}\) is governed by
\begin{equation}
\ddot{\bm{T}}_s \;=\;
 -A_s(t)\,\bm{\omega}
 \;-\; B_s(t)\,\dot{\bm{\omega}}
 \;-\; C_s(t)\,\dot{\bm{T}}_s
 \;-\; K_s(t)\,\bm{T}_s,
\label{eq:slosh}
\end{equation}
with each scalar coefficient time-varying as
\begin{equation}
f(t) \;=\; f_{\text{nom}} + f_{\text{amp}}\sin(\omega_f\,t),
\qquad f \in \{A_s,B_s,C_s,K_s\}.
\end{equation}
The constraint \(f_{\text{nom}} > |f_{\text{amp}}|\) keeps \(K_s, C_s\)
positive (negative stiffness or damping would be unphysical). \(\bm{T}_s\)
enters Euler's equation as an additional external torque on the body.

\textbf{No algebraic loop.} \(\dot{\bm{\omega}}\) depends on \(\bm{T}_s\) (a
\emph{state}, not on \(\ddot{\bm{T}}_s\)), and \(\ddot{\bm{T}}_s\) consumes
the just-computed \(\dot{\bm{\omega}}\). The RHS is fully explicit.

\subsection{Coupled Equations of Motion}

The complete state is
\(\bm{x} = [\bm{q};\bm{\omega};\bm{h}_w;\bm{T}_s;\dot{\bm{T}}_s]
        \in \mathbb{R}^{13+n_w}\),
governed by
\begin{align}
\dot{\bm{q}}        &= \tfrac{1}{2}\,\bm{\Omega}(\bm{\omega})\,\bm{q}, \\
\bm{J}\,\dot{\bm{\omega}} &= \bm{u}_{\text{body}}
                            - \bm{\omega}\times(\bm{J}\bm{\omega}+\bm{W}\bm{h}_w), \\
\dot{\bm{h}}_w      &= \bm{\tau}_w, \\
\ddot{\bm{T}}_s     &= -A_s\bm{\omega} - B_s\dot{\bm{\omega}}
                       - C_s\dot{\bm{T}}_s - K_s\bm{T}_s,
\end{align}
with
\(\bm{u}_{\text{body}} = -\bm{W}\bm{\tau}_w + \bm{m}\times\bm{B} + \bm{T}_s\).


% ---------------------------------------------------------------------
\section{Code Architecture Evolution}
% ---------------------------------------------------------------------

The simulation was built incrementally; each step is reproducible from
the conversation log.

\begin{enumerate}\setlength{\itemsep}{4pt}
  \item \textbf{Baseline rigid body + PD controller.}
        A single-file implementation of Euler's equation with a quaternion
        PD law and a hand-written fourth-order Runge--Kutta integrator.
  \item \textbf{Euler-angle initial conditions.}
        Helper \texttt{euler321\_to\_quat()} converts a roll-pitch-yaw
        triple to a scalar-last quaternion (M\&C App.~B.96).
  \item \textbf{Full 3$\times$3 inertia tensor.}
        Replaced the diagonal default with off-diagonal products of inertia
        and added symmetry / positive-definiteness assertions.
  \item \textbf{Reaction-wheel dynamics.}
        Added the four-wheel NASA pyramid configuration matrix \(\bm{W}\),
        pseudoinverse allocation, and the wheel-momentum gyroscopic term in
        Euler's equation.
  \item \textbf{Magnetic momentum unloading.}
        Added a toy LEO B-field, the cross-product unloading law, and per-
        axis dipole saturation. Also introduced PD torque saturation
        \(u_{\max}\).
  \item \textbf{Propellant slosh disturbance.}
        Added the second-order slosh ODE driven by \((\bm{\omega},
        \dot{\bm{\omega}})\), with time-varying scalar coefficients
        (DC~+~sine).
  \item \textbf{Variable rename.}
        Renamed \(\Gamma_s \to T_s\) throughout the code and plot labels
        for consistency with the user's preferred notation.
  \item \textbf{Toggles.}
        Added independent on/off switches for slosh, magnetorquer, and
        reaction wheels. With wheels disabled the PD output drives the body
        directly as an ideal torque actuator.
  \item \textbf{Migration to \texttt{solve\_ivp}.}
        Replaced the hand-rolled RK4 loop with
        \texttt{scipy.integrate.solve\_ivp} (adaptive RK45). Quaternion
        re-normalisation moved from per-step to post-processing.
  \item \textbf{Merge and clean-up.}
        Consolidated the two intermediate files into a single
        \texttt{rigid\_body\_pd.py}. Removed all RK4-specific code paths.
  \item \textbf{Dataclass + presets refactor (v2).}
        Centralised every spacecraft parameter, gain, initial condition,
        time setting, and toggle into a single \texttt{SimConfig} dataclass
        at the top of the file. Added two named preset factories,
        \texttt{mc\_example\_71()} and \texttt{mc\_example\_72()}, matching
        the worked examples in M\&C Sec.~7.
\end{enumerate}


% ---------------------------------------------------------------------
\section{Final Architecture (v2)}
% ---------------------------------------------------------------------

\subsection{The \texttt{SimConfig} Dataclass}

All scenario knobs live in one place:

\begin{lstlisting}
@dataclass
class SimConfig:
    """All satellite parameters, gains, ICs, time, and toggles."""
    # Inertia (M&C Eq. 3.13), kg.m^2
    Ixx: float = 6400.0;  Iyy: float = 4730.0;  Izz: float = 8160.0
    Ixy: float = 76.4;    Ixz: float = 25.6;    Iyz: float = 40.0
    # Wheel geometry
    wheel_tilt_deg: float = 54.7356      # 4-wheel NASA pyramid
    # PD gains and saturation
    Kp: float = 10.0;   Kd: float = 150.0;   u_max: float = 1000.0
    # Magnetorquer
    K_dump: float = 1.0;  m_max: float = 500.0
    # Slosh:  f(t) = nom + amp*sin(omega*t)
    slosh_params: dict = field(default_factory=lambda: {
        'A': {'nom': 0.005, 'amp': 0.0, 'omega': 0.05},
        'B': {'nom': 0.020, 'amp': 0.0, 'omega': 0.05},
        'C': {'nom': 0.023, 'amp': 0.0, 'omega': 0.05},
        'K': {'nom': 0.008, 'amp': 0.0, 'omega': 0.05},
    })
    # Initial conditions (scalar-last)
    q0:    np.ndarray = field(default_factory=lambda:
        (math.sqrt(2)/2) * np.array([1.0, 0.0, 0.0, 1.0]))
    w0:    np.ndarray = field(default_factory=lambda:
        np.array([0.01, 0.01, 0.01]))
    q_des: np.ndarray = field(default_factory=lambda:
        np.array([0.0, 0.0, 0.0, 1.0]))
    Ts0:   np.ndarray = field(default_factory=lambda: np.zeros(3))
    Tsd0:  np.ndarray = field(default_factory=lambda: np.zeros(3))
    # Time
    t_end: float = 20*60.0;  dt_eval: float = 0.05
    # Toggles
    enable_slosh:        bool = False
    enable_magnetorquer: bool = False
    enable_wheels:       bool = True
    # Solver
    method: str = 'RK45';  rtol: float = 1e-8
    atol: float = 1e-10;   max_step: float = 0.5
\end{lstlisting}

\subsection{Named Presets}

\begin{lstlisting}
def mc_example_71() -> SimConfig:
    """Large spacecraft, diagonal J, fast initial tumble."""
    cfg = SimConfig()
    cfg.Ixx, cfg.Iyy, cfg.Izz = 10000.0, 9000.0, 12000.0
    cfg.Ixy = cfg.Ixz = cfg.Iyz = 0.0
    cfg.Kp, cfg.Kd = 50.0, 500.0
    cfg.u_max = 100.0
    q = np.array([0.6853, 0.6953, 0.1531, 0.1531])
    cfg.q0 = q / np.linalg.norm(q)
    cfg.w0 = np.array([0.53, 0.53, 0.053])
    cfg.t_end = 300.0
    return cfg

def mc_example_72() -> SimConfig:
    """Smaller spacecraft with products of inertia; gentle IC."""
    return SimConfig()    # SimConfig() defaults already match Example 7.2
\end{lstlisting}

\subsection{Usage Pattern}

\begin{lstlisting}
cfg = mc_example_72()
cfg.enable_slosh = True            # in-place tweak
cfg.Kp, cfg.Kd = 15.0, 200.0
results = simulate(cfg)            # one entry point
plot_results(*results)
plot_euler_and_rates(results[0], results[1], results[3])
plt.show()
\end{lstlisting}


% ---------------------------------------------------------------------
\section{Validation: Automated Smoke Test}
% ---------------------------------------------------------------------

A separate file \texttt{test\_rigid\_body\_pd\_v2.py} exercises the
simulation under every toggle combination using the non-interactive
\texttt{Agg} matplotlib backend. All eight test cases pass.

\begin{center}
\begin{tabular}{lc}
\toprule
\textbf{Test} & \textbf{Result} \\
\midrule
Euler $\leftrightarrow$ quaternion roundtrip (4 angle triples) & Pass \\
Example 7.1 preset, wheels only                                 & Pass \\
Example 7.2 preset, wheels only                                 & Pass \\
Example 7.2 full stack (slosh + mag + wheels)                   & Pass \\
Example 7.2 with wheels OFF (ideal actuator)                    & Pass \\
Example 7.2 with all subsystems OFF                             & Pass \\
Example 7.2 with overridden gains (in-place mutation)           & Pass \\
Plot functions (both figures)                                   & Pass \\
\bottomrule
\end{tabular}
\end{center}

\noindent Per-trajectory invariants checked at every output sample:
\begin{itemize}\setlength{\itemsep}{1pt}
  \item Output tuple length = 11; all array shapes match the expected
        \((N,4)\), \((N,3)\), \((N,n_w)\).
  \item No NaNs in any state trajectory.
  \item Unit-norm quaternion preserved to machine precision
        (\(\max\bigl|\,\|\bm{q}\|-1\,\bigr| = 2.22\times10^{-16}\)).
  \item \(|\bm{\omega}|\) monotonically non-increasing in every settling case.
  \item When a toggle is OFF, the corresponding telemetry stays at
        exactly zero (\(\bm{h}_w\), \(\bm{\tau}_w\), \(\bm{m}\),
        \(\bm{\tau}_{\text{mag}}\), \(\bm{T}_s\)).
\end{itemize}


% ---------------------------------------------------------------------
\section{Companion LaTeX Physics Summary}
% ---------------------------------------------------------------------

A separate file \texttt{attitude\_dynamics\_model.tex} was produced for
direct import into the project's Overleaf write-up. It contains a
self-contained section with every equation above plus extended
discussion, fully cross-referenced via \texttt{\textbackslash label}.
The \texttt{\textbackslash cite\{markley2014\}} entry should resolve to
Markley and Crassidis using the BibTeX stanza included at the top of that
file.


% ---------------------------------------------------------------------
\section{Proposed Extension: NARX Slosh Surrogate}
% ---------------------------------------------------------------------

A design (no code) was discussed for replacing the analytical slosh ODE
\eqref{eq:slosh} with a learned Nonlinear Autoregressive Network with
eXogenous Inputs (NARX). The map to learn is
\begin{equation}
\hat{T}_s(k) \;=\; f_\theta\bigl(
  T_s(k{-}1),\dots,T_s(k{-}n_y),\;
  \omega(k),\dots,\omega(k{-}n_u),\;
  \dot\omega(k),\dots,\dot\omega(k{-}n_u)
\bigr),
\end{equation}
optionally extended to a \emph{parametric} surrogate by also feeding the
four slosh coefficients \((A_s, B_s, C_s, K_s)\). The plan, in order:

\begin{enumerate}\setlength{\itemsep}{2pt}
  \item Add a helper \texttt{reconstruct\_w\_dot()} that recomputes
        \(\dot{\bm{\omega}}\) exactly from the saved \texttt{simulate()}
        outputs (no finite differences).
  \item Implement \texttt{sample\_random\_config(rng)} that draws
        \texttt{SimConfig} instances with randomised initial conditions,
        gains, inertia perturbations, and (optionally) slosh coefficients.
  \item Generate 100--500 trajectories, each $\sim$600\,s, downsampled to
        \(\Delta t \approx 0.5\,\mathrm{s}\) and stacked across all three
        body axes (the slosh dynamics are axis-decoupled, so one scalar
        NARX trained on stacked per-axis data is the right architecture).
  \item Build sliding-window features with lag orders
        \(n_y \approx n_u \approx 10\)--\(20\) samples.
  \item Train a small MLP (e.g. 64-64 \texttt{tanh}) with
        \texttt{sklearn.neural\_network.MLPRegressor} as a baseline, then
        promote to PyTorch if needed.
  \item Validate in two modes:
        \emph{one-step prediction} (teacher-forced lags) for training-loss
        sanity, and \emph{free-running rollout} (model feeds its own past
        predictions) for honest dynamic accuracy.
  \item Split by \emph{whole trajectory} (70/15/15) to avoid leakage from
        adjacent correlated samples within one simulation.
  \item Optional integration: add \texttt{cfg.use\_narx\_slosh = True}
        and \texttt{cfg.narx\_model = trained\_model} to the dataclass,
        wire it into \texttt{dynamics\_rhs}, and compare closed-loop
        traces against the analytical slosh ground truth on held-out
        initial conditions.
\end{enumerate}


% ---------------------------------------------------------------------
\section{Deliverables}
% ---------------------------------------------------------------------

\begin{center}
\begin{tabular}{ll}
\toprule
\textbf{File} & \textbf{Purpose} \\
\midrule
\texttt{rigid\_body\_pd.py}            & Merged simulation, \texttt{solve\_ivp}-based \\
\texttt{rigid\_body\_pd\_v2.py}        & Dataclass + presets refactor (recommended) \\
\texttt{test\_rigid\_body\_pd\_v2.py}  & Smoke test (8 cases, all pass) \\
\texttt{attitude\_dynamics\_model.tex} & Physics section for Overleaf \\
\texttt{session\_summary.tex}          & This document \\
\bottomrule
\end{tabular}
\end{center}


% ---------------------------------------------------------------------
\section{Conventions Used Throughout}
% ---------------------------------------------------------------------

\begin{itemize}\setlength{\itemsep}{1pt}
  \item \textbf{Quaternion:} scalar-last, \(\bm{q}=[q_1, q_2, q_3, q_4]\)
        with \(q_4\) the scalar part.
  \item \textbf{Euler sequence:} 3-2-1 (Z-Y-X, yaw-pitch-roll) for human
        I/O; quaternions everywhere internally.
  \item \textbf{Inertia sign:} off-diagonal entries are stored as positive
        products of inertia; the negative signs live in the \(\bm{J}\)
        builder (M\&C Eq.~3.13).
  \item \textbf{Frames:} body frame throughout; \(\bm{B}\) and
        \(\bm{h}_w^{\,b}\) are body-frame projections.
  \item \textbf{Slosh torque symbol:} \(\bm{T}_s\) (code: \texttt{Ts}),
        with rate \(\dot{\bm{T}}_s\) (\texttt{Ts\_dot}). Renamed from the
        original \(\Gamma_s\).
\end{itemize}


% ---------------------------------------------------------------------
\begin{thebibliography}{1}
% ---------------------------------------------------------------------

\bibitem{markley2014}
F.~L.~Markley and J.~L.~Crassidis,
\emph{Fundamentals of Spacecraft Attitude Determination and Control},
Space Technology Library, Vol.~33,
Springer, New York, NY, 2014.
\href{https://doi.org/10.1007/978-1-4939-0802-8}{doi:10.1007/978-1-4939-0802-8}.

\end{thebibliography}

\end{document}
 into your existing project.
%
% BIBLIOGRAPHY
%   The single \cite{markley2014} in this document needs the following
%   entry in your .bib file:
%
%   @book{markley2014,
%     author    = {Markley, F. Landis and Crassidis, John L.},
%     title     = {Fundamentals of Spacecraft Attitude Determination and Control},
%     publisher = {Springer},
%     year      = {2014},
%     series    = {Space Technology Library},
%     volume    = {33},
%     address   = {New York, NY},
%     doi       = {10.1007/978-1-4939-0802-8}
%   }
%
%   If you don't use BibTeX, find-and-replace \cite{markley2014} with
%   "Markley \& Crassidis (2014)".
% =====================================================================

\documentclass[11pt]{article}

\usepackage[margin=1in]{geometry}
\usepackage{amsmath,amssymb,bm}
\usepackage{listings}
\usepackage{xcolor}
\usepackage{booktabs}
\usepackage{enumitem}
\usepackage{hyperref}
\usepackage{tcolorbox}

\hypersetup{colorlinks=true, linkcolor=blue!60!black, urlcolor=blue!60!black}

\lstdefinestyle{pycode}{
    language=Python,
    basicstyle=\ttfamily\footnotesize,
    keywordstyle=\color{blue!70!black}\bfseries,
    stringstyle=\color{teal!70!black},
    commentstyle=\color{gray}\itshape,
    showstringspaces=false,
    breaklines=true,
    frame=single,
    rulecolor=\color{black!20},
    tabsize=4,
    columns=flexible,
    aboveskip=8pt,
    belowskip=8pt,
}
\lstset{style=pycode}

\title{Spacecraft Attitude Dynamics and Control Simulation\\[3pt]
       \large Session Build-Up Summary}
\author{}
\date{\today}


\begin{document}
\maketitle


% ---------------------------------------------------------------------
\section{Project Goal}
% ---------------------------------------------------------------------

Build, incrementally, a self-contained Python simulation of spacecraft
attitude dynamics and control following Markley and
Crassidis~\cite{markley2014}. Each step in the session added a piece of
physics or refactored the code for clarity. The final product is two
deliverables --- a working simulation in Python and a companion
LaTeX summary of the underlying physics --- plus an automated smoke test.


% ---------------------------------------------------------------------
\section{System Overview}
% ---------------------------------------------------------------------

The final simulation models a rigid spacecraft with the following
components:

\begin{itemize}\setlength{\itemsep}{2pt}
  \item Quaternion-feedback PD attitude controller (scalar-last convention).
  \item Four reaction wheels in a NASA-standard pyramid configuration.
  \item Magnetic torquers performing wheel momentum unloading via a
        cross-product (``\(B\)-dot--style'') control law.
  \item A propellant-slosh disturbance modelled as a forced second-order
        ODE on a 3-D torque state.
  \item A 3-2-1 (Z-Y-X) Euler-angle interface for human-readable initial
        conditions and post-processing.
\end{itemize}

\noindent The integrator is \texttt{scipy.integrate.solve\_ivp} (adaptive
RK45 by default), and the closed-loop is treated as continuous-time --- the
PD, magnetorquer, and wheel-allocation laws are evaluated at every solver
step rather than held with zero-order hold.


% ---------------------------------------------------------------------
\section{Modelling Components}
% ---------------------------------------------------------------------

\subsection{Quaternion Kinematics (Scalar-Last)}

Using \(\bm{q}=[\bm{q}_v;\,q_4]\), the body-to-inertial quaternion satisfies
\begin{equation}
\dot{\bm{q}} = \tfrac{1}{2}\,\bm{\Omega}(\bm{\omega})\,\bm{q}
\;=\;
\begin{bmatrix}
 \tfrac{1}{2}\bigl(q_4\,\bm{\omega} - \bm{\omega}\times\bm{q}_v\bigr) \\
-\tfrac{1}{2}\,\bm{\omega}^{\!\top}\bm{q}_v
\end{bmatrix},
\tag{M\&C Eq.~2.88}
\end{equation}
and the quaternion product follows
\(\bm{A}(\bm{q}\otimes\bm{p}) = \bm{A}(\bm{q})\bm{A}(\bm{p})\)
(M\&C~Eq.~2.82b).

\subsection{Rigid-Body Dynamics with Reaction Wheels}

Euler's equation with internal angular momentum
\(\bm{h}_w^{\,b} = \bm{W}\bm{h}_w\) is
\begin{equation}
\bm{J}\,\dot{\bm{\omega}}
\;=\;
\bm{u}_{\text{body}}
\;-\;
\bm{\omega}\times(\bm{J}\bm{\omega} + \bm{h}_w^{\,b}),
\tag{M\&C Sec.~7.2}
\end{equation}
where the inertia tensor uses the sign convention of M\&C~Eq.~3.13.
Wheel torque commands integrate the wheel-momentum law
\(\dot{\bm{h}}_w = \bm{\tau}_w\).

\subsection{Reaction-Wheel Allocation}

A four-wheel NASA pyramid with cant angle \(\beta = \arccos(1/\sqrt{3})
\approx 54.74^\circ\) gives equal projection of each wheel onto the body
axes. A desired body torque is allocated by the minimum-norm pseudoinverse:
\begin{equation}
\bm{\tau}_w \;=\; -\bm{W}^{+}\,\bm{u}_{\text{cmd}},
\qquad
\bm{W}^{+} = \bm{W}^{\!\top}(\bm{W}\bm{W}^{\!\top})^{-1}.
\end{equation}

\subsection{Quaternion-Feedback PD Controller}

With error quaternion \(\delta\bm{q} = \bm{q}_{\text{des}}^{-1}\otimes\bm{q}\),
the controller (M\&C~Sec.~7.4) is
\begin{equation}
\bm{u}_{\text{cmd}} \;=\;
\operatorname{sat}_{u_{\max}}\!\Bigl(
 -K_p\,\operatorname{sign}(\delta q_4)\,\delta\bm{q}_v
 \;-\; K_d\,\bm{\omega}
\Bigr).
\end{equation}
The \(\operatorname{sign}(\delta q_4)\) factor selects the shortest rotation
and prevents the well-known \emph{unwinding} phenomenon. Each component is
clipped to \([-u_{\max},+u_{\max}]\).

\subsection{Magnetic Momentum Unloading}

The cross-product (``\(B\)-dot--style'') law from M\&C~Sec.~7.5
\begin{equation}
\bm{m} \;=\; \operatorname{sat}_{m_{\max}}\!\biggl(
\frac{K_{\text{dump}}}{\|\bm{B}\|^{2}}\,
\bm{h}_w^{\,b}\times\bm{B}\biggr),
\quad
\bm{\tau}_{\text{mag}} \;=\; \bm{m}\times\bm{B},
\end{equation}
applies a damping torque to the component of \(\bm{h}_w^{\,b}\)
perpendicular to \(\bm{B}\). Orbital variation of \(\bm{B}\) sweeps the
unloadable subspace through all three body axes over one orbit. The toy
geomagnetic field is an inclined rotating dipole of magnitude
\(\sim 30~\mu\mathrm{T}\) at 90-min LEO period.

\subsection{Propellant Slosh Model}

Slosh torque \(\bm{T}_s \in \mathbb{R}^{3}\) is governed by
\begin{equation}
\ddot{\bm{T}}_s \;=\;
 -A_s(t)\,\bm{\omega}
 \;-\; B_s(t)\,\dot{\bm{\omega}}
 \;-\; C_s(t)\,\dot{\bm{T}}_s
 \;-\; K_s(t)\,\bm{T}_s,
\label{eq:slosh}
\end{equation}
with each scalar coefficient time-varying as
\begin{equation}
f(t) \;=\; f_{\text{nom}} + f_{\text{amp}}\sin(\omega_f\,t),
\qquad f \in \{A_s,B_s,C_s,K_s\}.
\end{equation}
The constraint \(f_{\text{nom}} > |f_{\text{amp}}|\) keeps \(K_s, C_s\)
positive (negative stiffness or damping would be unphysical). \(\bm{T}_s\)
enters Euler's equation as an additional external torque on the body.

\textbf{No algebraic loop.} \(\dot{\bm{\omega}}\) depends on \(\bm{T}_s\) (a
\emph{state}, not on \(\ddot{\bm{T}}_s\)), and \(\ddot{\bm{T}}_s\) consumes
the just-computed \(\dot{\bm{\omega}}\). The RHS is fully explicit.

\subsection{Coupled Equations of Motion}

The complete state is
\(\bm{x} = [\bm{q};\bm{\omega};\bm{h}_w;\bm{T}_s;\dot{\bm{T}}_s]
        \in \mathbb{R}^{13+n_w}\),
governed by
\begin{align}
\dot{\bm{q}}        &= \tfrac{1}{2}\,\bm{\Omega}(\bm{\omega})\,\bm{q}, \\
\bm{J}\,\dot{\bm{\omega}} &= \bm{u}_{\text{body}}
                            - \bm{\omega}\times(\bm{J}\bm{\omega}+\bm{W}\bm{h}_w), \\
\dot{\bm{h}}_w      &= \bm{\tau}_w, \\
\ddot{\bm{T}}_s     &= -A_s\bm{\omega} - B_s\dot{\bm{\omega}}
                       - C_s\dot{\bm{T}}_s - K_s\bm{T}_s,
\end{align}
with
\(\bm{u}_{\text{body}} = -\bm{W}\bm{\tau}_w + \bm{m}\times\bm{B} + \bm{T}_s\).


% ---------------------------------------------------------------------
\section{Code Architecture Evolution}
% ---------------------------------------------------------------------

The simulation was built incrementally; each step is reproducible from
the conversation log.

\begin{enumerate}\setlength{\itemsep}{4pt}
  \item \textbf{Baseline rigid body + PD controller.}
        A single-file implementation of Euler's equation with a quaternion
        PD law and a hand-written fourth-order Runge--Kutta integrator.
  \item \textbf{Euler-angle initial conditions.}
        Helper \texttt{euler321\_to\_quat()} converts a roll-pitch-yaw
        triple to a scalar-last quaternion (M\&C App.~B.96).
  \item \textbf{Full 3$\times$3 inertia tensor.}
        Replaced the diagonal default with off-diagonal products of inertia
        and added symmetry / positive-definiteness assertions.
  \item \textbf{Reaction-wheel dynamics.}
        Added the four-wheel NASA pyramid configuration matrix \(\bm{W}\),
        pseudoinverse allocation, and the wheel-momentum gyroscopic term in
        Euler's equation.
  \item \textbf{Magnetic momentum unloading.}
        Added a toy LEO B-field, the cross-product unloading law, and per-
        axis dipole saturation. Also introduced PD torque saturation
        \(u_{\max}\).
  \item \textbf{Propellant slosh disturbance.}
        Added the second-order slosh ODE driven by \((\bm{\omega},
        \dot{\bm{\omega}})\), with time-varying scalar coefficients
        (DC~+~sine).
  \item \textbf{Variable rename.}
        Renamed \(\Gamma_s \to T_s\) throughout the code and plot labels
        for consistency with the user's preferred notation.
  \item \textbf{Toggles.}
        Added independent on/off switches for slosh, magnetorquer, and
        reaction wheels. With wheels disabled the PD output drives the body
        directly as an ideal torque actuator.
  \item \textbf{Migration to \texttt{solve\_ivp}.}
        Replaced the hand-rolled RK4 loop with
        \texttt{scipy.integrate.solve\_ivp} (adaptive RK45). Quaternion
        re-normalisation moved from per-step to post-processing.
  \item \textbf{Merge and clean-up.}
        Consolidated the two intermediate files into a single
        \texttt{rigid\_body\_pd.py}. Removed all RK4-specific code paths.
  \item \textbf{Dataclass + presets refactor (v2).}
        Centralised every spacecraft parameter, gain, initial condition,
        time setting, and toggle into a single \texttt{SimConfig} dataclass
        at the top of the file. Added two named preset factories,
        \texttt{mc\_example\_71()} and \texttt{mc\_example\_72()}, matching
        the worked examples in M\&C Sec.~7.
\end{enumerate}


% ---------------------------------------------------------------------
\section{Final Architecture (v2)}
% ---------------------------------------------------------------------

\subsection{The \texttt{SimConfig} Dataclass}

All scenario knobs live in one place:

\begin{lstlisting}
@dataclass
class SimConfig:
    """All satellite parameters, gains, ICs, time, and toggles."""
    # Inertia (M&C Eq. 3.13), kg.m^2
    Ixx: float = 6400.0;  Iyy: float = 4730.0;  Izz: float = 8160.0
    Ixy: float = 76.4;    Ixz: float = 25.6;    Iyz: float = 40.0
    # Wheel geometry
    wheel_tilt_deg: float = 54.7356      # 4-wheel NASA pyramid
    # PD gains and saturation
    Kp: float = 10.0;   Kd: float = 150.0;   u_max: float = 1000.0
    # Magnetorquer
    K_dump: float = 1.0;  m_max: float = 500.0
    # Slosh:  f(t) = nom + amp*sin(omega*t)
    slosh_params: dict = field(default_factory=lambda: {
        'A': {'nom': 0.005, 'amp': 0.0, 'omega': 0.05},
        'B': {'nom': 0.020, 'amp': 0.0, 'omega': 0.05},
        'C': {'nom': 0.023, 'amp': 0.0, 'omega': 0.05},
        'K': {'nom': 0.008, 'amp': 0.0, 'omega': 0.05},
    })
    # Initial conditions (scalar-last)
    q0:    np.ndarray = field(default_factory=lambda:
        (math.sqrt(2)/2) * np.array([1.0, 0.0, 0.0, 1.0]))
    w0:    np.ndarray = field(default_factory=lambda:
        np.array([0.01, 0.01, 0.01]))
    q_des: np.ndarray = field(default_factory=lambda:
        np.array([0.0, 0.0, 0.0, 1.0]))
    Ts0:   np.ndarray = field(default_factory=lambda: np.zeros(3))
    Tsd0:  np.ndarray = field(default_factory=lambda: np.zeros(3))
    # Time
    t_end: float = 20*60.0;  dt_eval: float = 0.05
    # Toggles
    enable_slosh:        bool = False
    enable_magnetorquer: bool = False
    enable_wheels:       bool = True
    # Solver
    method: str = 'RK45';  rtol: float = 1e-8
    atol: float = 1e-10;   max_step: float = 0.5
\end{lstlisting}

\subsection{Named Presets}

\begin{lstlisting}
def mc_example_71() -> SimConfig:
    """Large spacecraft, diagonal J, fast initial tumble."""
    cfg = SimConfig()
    cfg.Ixx, cfg.Iyy, cfg.Izz = 10000.0, 9000.0, 12000.0
    cfg.Ixy = cfg.Ixz = cfg.Iyz = 0.0
    cfg.Kp, cfg.Kd = 50.0, 500.0
    cfg.u_max = 100.0
    q = np.array([0.6853, 0.6953, 0.1531, 0.1531])
    cfg.q0 = q / np.linalg.norm(q)
    cfg.w0 = np.array([0.53, 0.53, 0.053])
    cfg.t_end = 300.0
    return cfg

def mc_example_72() -> SimConfig:
    """Smaller spacecraft with products of inertia; gentle IC."""
    return SimConfig()    # SimConfig() defaults already match Example 7.2
\end{lstlisting}

\subsection{Usage Pattern}

\begin{lstlisting}
cfg = mc_example_72()
cfg.enable_slosh = True            # in-place tweak
cfg.Kp, cfg.Kd = 15.0, 200.0
results = simulate(cfg)            # one entry point
plot_results(*results)
plot_euler_and_rates(results[0], results[1], results[3])
plt.show()
\end{lstlisting}


% ---------------------------------------------------------------------
\section{Validation: Automated Smoke Test}
% ---------------------------------------------------------------------

A separate file \texttt{test\_rigid\_body\_pd\_v2.py} exercises the
simulation under every toggle combination using the non-interactive
\texttt{Agg} matplotlib backend. All eight test cases pass.

\begin{center}
\begin{tabular}{lc}
\toprule
\textbf{Test} & \textbf{Result} \\
\midrule
Euler $\leftrightarrow$ quaternion roundtrip (4 angle triples) & Pass \\
Example 7.1 preset, wheels only                                 & Pass \\
Example 7.2 preset, wheels only                                 & Pass \\
Example 7.2 full stack (slosh + mag + wheels)                   & Pass \\
Example 7.2 with wheels OFF (ideal actuator)                    & Pass \\
Example 7.2 with all subsystems OFF                             & Pass \\
Example 7.2 with overridden gains (in-place mutation)           & Pass \\
Plot functions (both figures)                                   & Pass \\
\bottomrule
\end{tabular}
\end{center}

\noindent Per-trajectory invariants checked at every output sample:
\begin{itemize}\setlength{\itemsep}{1pt}
  \item Output tuple length = 11; all array shapes match the expected
        \((N,4)\), \((N,3)\), \((N,n_w)\).
  \item No NaNs in any state trajectory.
  \item Unit-norm quaternion preserved to machine precision
        (\(\max\bigl|\,\|\bm{q}\|-1\,\bigr| = 2.22\times10^{-16}\)).
  \item \(|\bm{\omega}|\) monotonically non-increasing in every settling case.
  \item When a toggle is OFF, the corresponding telemetry stays at
        exactly zero (\(\bm{h}_w\), \(\bm{\tau}_w\), \(\bm{m}\),
        \(\bm{\tau}_{\text{mag}}\), \(\bm{T}_s\)).
\end{itemize}


% ---------------------------------------------------------------------
\section{Companion LaTeX Physics Summary}
% ---------------------------------------------------------------------

A separate file \texttt{attitude\_dynamics\_model.tex} was produced for
direct import into the project's Overleaf write-up. It contains a
self-contained section with every equation above plus extended
discussion, fully cross-referenced via \texttt{\textbackslash label}.
The \texttt{\textbackslash cite\{markley2014\}} entry should resolve to
Markley and Crassidis using the BibTeX stanza included at the top of that
file.


% ---------------------------------------------------------------------
\section{Proposed Extension: NARX Slosh Surrogate}
% ---------------------------------------------------------------------

A design (no code) was discussed for replacing the analytical slosh ODE
\eqref{eq:slosh} with a learned Nonlinear Autoregressive Network with
eXogenous Inputs (NARX). The map to learn is
\begin{equation}
\hat{T}_s(k) \;=\; f_\theta\bigl(
  T_s(k{-}1),\dots,T_s(k{-}n_y),\;
  \omega(k),\dots,\omega(k{-}n_u),\;
  \dot\omega(k),\dots,\dot\omega(k{-}n_u)
\bigr),
\end{equation}
optionally extended to a \emph{parametric} surrogate by also feeding the
four slosh coefficients \((A_s, B_s, C_s, K_s)\). The plan, in order:

\begin{enumerate}\setlength{\itemsep}{2pt}
  \item Add a helper \texttt{reconstruct\_w\_dot()} that recomputes
        \(\dot{\bm{\omega}}\) exactly from the saved \texttt{simulate()}
        outputs (no finite differences).
  \item Implement \texttt{sample\_random\_config(rng)} that draws
        \texttt{SimConfig} instances with randomised initial conditions,
        gains, inertia perturbations, and (optionally) slosh coefficients.
  \item Generate 100--500 trajectories, each $\sim$600\,s, downsampled to
        \(\Delta t \approx 0.5\,\mathrm{s}\) and stacked across all three
        body axes (the slosh dynamics are axis-decoupled, so one scalar
        NARX trained on stacked per-axis data is the right architecture).
  \item Build sliding-window features with lag orders
        \(n_y \approx n_u \approx 10\)--\(20\) samples.
  \item Train a small MLP (e.g. 64-64 \texttt{tanh}) with
        \texttt{sklearn.neural\_network.MLPRegressor} as a baseline, then
        promote to PyTorch if needed.
  \item Validate in two modes:
        \emph{one-step prediction} (teacher-forced lags) for training-loss
        sanity, and \emph{free-running rollout} (model feeds its own past
        predictions) for honest dynamic accuracy.
  \item Split by \emph{whole trajectory} (70/15/15) to avoid leakage from
        adjacent correlated samples within one simulation.
  \item Optional integration: add \texttt{cfg.use\_narx\_slosh = True}
        and \texttt{cfg.narx\_model = trained\_model} to the dataclass,
        wire it into \texttt{dynamics\_rhs}, and compare closed-loop
        traces against the analytical slosh ground truth on held-out
        initial conditions.
\end{enumerate}


% ---------------------------------------------------------------------
\section{Deliverables}
% ---------------------------------------------------------------------

\begin{center}
\begin{tabular}{ll}
\toprule
\textbf{File} & \textbf{Purpose} \\
\midrule
\texttt{rigid\_body\_pd.py}            & Merged simulation, \texttt{solve\_ivp}-based \\
\texttt{rigid\_body\_pd\_v2.py}        & Dataclass + presets refactor (recommended) \\
\texttt{test\_rigid\_body\_pd\_v2.py}  & Smoke test (8 cases, all pass) \\
\texttt{attitude\_dynamics\_model.tex} & Physics section for Overleaf \\
\texttt{session\_summary.tex}          & This document \\
\bottomrule
\end{tabular}
\end{center}


% ---------------------------------------------------------------------
\section{Conventions Used Throughout}
% ---------------------------------------------------------------------

\begin{itemize}\setlength{\itemsep}{1pt}
  \item \textbf{Quaternion:} scalar-last, \(\bm{q}=[q_1, q_2, q_3, q_4]\)
        with \(q_4\) the scalar part.
  \item \textbf{Euler sequence:} 3-2-1 (Z-Y-X, yaw-pitch-roll) for human
        I/O; quaternions everywhere internally.
  \item \textbf{Inertia sign:} off-diagonal entries are stored as positive
        products of inertia; the negative signs live in the \(\bm{J}\)
        builder (M\&C Eq.~3.13).
  \item \textbf{Frames:} body frame throughout; \(\bm{B}\) and
        \(\bm{h}_w^{\,b}\) are body-frame projections.
  \item \textbf{Slosh torque symbol:} \(\bm{T}_s\) (code: \texttt{Ts}),
        with rate \(\dot{\bm{T}}_s\) (\texttt{Ts\_dot}). Renamed from the
        original \(\Gamma_s\).
\end{itemize}


% ---------------------------------------------------------------------
\begin{thebibliography}{1}
% ---------------------------------------------------------------------

\bibitem{markley2014}
F.~L.~Markley and J.~L.~Crassidis,
\emph{Fundamentals of Spacecraft Attitude Determination and Control},
Space Technology Library, Vol.~33,
Springer, New York, NY, 2014.
\href{https://doi.org/10.1007/978-1-4939-0802-8}{doi:10.1007/978-1-4939-0802-8}.

\end{thebibliography}

\end{document}
 into your existing project.
%
% BIBLIOGRAPHY
%   The single \cite{markley2014} in this document needs the following
%   entry in your .bib file:
%
%   @book{markley2014,
%     author    = {Markley, F. Landis and Crassidis, John L.},
%     title     = {Fundamentals of Spacecraft Attitude Determination and Control},
%     publisher = {Springer},
%     year      = {2014},
%     series    = {Space Technology Library},
%     volume    = {33},
%     address   = {New York, NY},
%     doi       = {10.1007/978-1-4939-0802-8}
%   }
%
%   If you don't use BibTeX, find-and-replace \cite{markley2014} with
%   "Markley \& Crassidis (2014)".
% =====================================================================

\documentclass[11pt]{article}

\usepackage[margin=1in]{geometry}
\usepackage{amsmath,amssymb,bm}
\usepackage{listings}
\usepackage{xcolor}
\usepackage{booktabs}
\usepackage{enumitem}
\usepackage{hyperref}
\usepackage{tcolorbox}

\hypersetup{colorlinks=true, linkcolor=blue!60!black, urlcolor=blue!60!black}

\lstdefinestyle{pycode}{
    language=Python,
    basicstyle=\ttfamily\footnotesize,
    keywordstyle=\color{blue!70!black}\bfseries,
    stringstyle=\color{teal!70!black},
    commentstyle=\color{gray}\itshape,
    showstringspaces=false,
    breaklines=true,
    frame=single,
    rulecolor=\color{black!20},
    tabsize=4,
    columns=flexible,
    aboveskip=8pt,
    belowskip=8pt,
}
\lstset{style=pycode}

\title{Spacecraft Attitude Dynamics and Control Simulation\\[3pt]
       \large Session Build-Up Summary}
\author{}
\date{\today}


\begin{document}
\maketitle


% ---------------------------------------------------------------------
\section{Project Goal}
% ---------------------------------------------------------------------

Build, incrementally, a self-contained Python simulation of spacecraft
attitude dynamics and control following Markley and
Crassidis~\cite{markley2014}. Each step in the session added a piece of
physics or refactored the code for clarity. The final product is two
deliverables --- a working simulation in Python and a companion
LaTeX summary of the underlying physics --- plus an automated smoke test.


% ---------------------------------------------------------------------
\section{System Overview}
% ---------------------------------------------------------------------

The final simulation models a rigid spacecraft with the following
components:

\begin{itemize}\setlength{\itemsep}{2pt}
  \item Quaternion-feedback PD attitude controller (scalar-last convention).
  \item Four reaction wheels in a NASA-standard pyramid configuration.
  \item Magnetic torquers performing wheel momentum unloading via a
        cross-product (``\(B\)-dot--style'') control law.
  \item A propellant-slosh disturbance modelled as a forced second-order
        ODE on a 3-D torque state.
  \item A 3-2-1 (Z-Y-X) Euler-angle interface for human-readable initial
        conditions and post-processing.
\end{itemize}

\noindent The integrator is \texttt{scipy.integrate.solve\_ivp} (adaptive
RK45 by default), and the closed-loop is treated as continuous-time --- the
PD, magnetorquer, and wheel-allocation laws are evaluated at every solver
step rather than held with zero-order hold.


% ---------------------------------------------------------------------
\section{Modelling Components}
% ---------------------------------------------------------------------

\subsection{Quaternion Kinematics (Scalar-Last)}

Using \(\bm{q}=[\bm{q}_v;\,q_4]\), the body-to-inertial quaternion satisfies
\begin{equation}
\dot{\bm{q}} = \tfrac{1}{2}\,\bm{\Omega}(\bm{\omega})\,\bm{q}
\;=\;
\begin{bmatrix}
 \tfrac{1}{2}\bigl(q_4\,\bm{\omega} - \bm{\omega}\times\bm{q}_v\bigr) \\
-\tfrac{1}{2}\,\bm{\omega}^{\!\top}\bm{q}_v
\end{bmatrix},
\tag{M\&C Eq.~2.88}
\end{equation}
and the quaternion product follows
\(\bm{A}(\bm{q}\otimes\bm{p}) = \bm{A}(\bm{q})\bm{A}(\bm{p})\)
(M\&C~Eq.~2.82b).

\subsection{Rigid-Body Dynamics with Reaction Wheels}

Euler's equation with internal angular momentum
\(\bm{h}_w^{\,b} = \bm{W}\bm{h}_w\) is
\begin{equation}
\bm{J}\,\dot{\bm{\omega}}
\;=\;
\bm{u}_{\text{body}}
\;-\;
\bm{\omega}\times(\bm{J}\bm{\omega} + \bm{h}_w^{\,b}),
\tag{M\&C Sec.~7.2}
\end{equation}
where the inertia tensor uses the sign convention of M\&C~Eq.~3.13.
Wheel torque commands integrate the wheel-momentum law
\(\dot{\bm{h}}_w = \bm{\tau}_w\).

\subsection{Reaction-Wheel Allocation}

A four-wheel NASA pyramid with cant angle \(\beta = \arccos(1/\sqrt{3})
\approx 54.74^\circ\) gives equal projection of each wheel onto the body
axes. A desired body torque is allocated by the minimum-norm pseudoinverse:
\begin{equation}
\bm{\tau}_w \;=\; -\bm{W}^{+}\,\bm{u}_{\text{cmd}},
\qquad
\bm{W}^{+} = \bm{W}^{\!\top}(\bm{W}\bm{W}^{\!\top})^{-1}.
\end{equation}

\subsection{Quaternion-Feedback PD Controller}

With error quaternion \(\delta\bm{q} = \bm{q}_{\text{des}}^{-1}\otimes\bm{q}\),
the controller (M\&C~Sec.~7.4) is
\begin{equation}
\bm{u}_{\text{cmd}} \;=\;
\operatorname{sat}_{u_{\max}}\!\Bigl(
 -K_p\,\operatorname{sign}(\delta q_4)\,\delta\bm{q}_v
 \;-\; K_d\,\bm{\omega}
\Bigr).
\end{equation}
The \(\operatorname{sign}(\delta q_4)\) factor selects the shortest rotation
and prevents the well-known \emph{unwinding} phenomenon. Each component is
clipped to \([-u_{\max},+u_{\max}]\).

\subsection{Magnetic Momentum Unloading}

The cross-product (``\(B\)-dot--style'') law from M\&C~Sec.~7.5
\begin{equation}
\bm{m} \;=\; \operatorname{sat}_{m_{\max}}\!\biggl(
\frac{K_{\text{dump}}}{\|\bm{B}\|^{2}}\,
\bm{h}_w^{\,b}\times\bm{B}\biggr),
\quad
\bm{\tau}_{\text{mag}} \;=\; \bm{m}\times\bm{B},
\end{equation}
applies a damping torque to the component of \(\bm{h}_w^{\,b}\)
perpendicular to \(\bm{B}\). Orbital variation of \(\bm{B}\) sweeps the
unloadable subspace through all three body axes over one orbit. The toy
geomagnetic field is an inclined rotating dipole of magnitude
\(\sim 30~\mu\mathrm{T}\) at 90-min LEO period.

\subsection{Propellant Slosh Model}

Slosh torque \(\bm{T}_s \in \mathbb{R}^{3}\) is governed by
\begin{equation}
\ddot{\bm{T}}_s \;=\;
 -A_s(t)\,\bm{\omega}
 \;-\; B_s(t)\,\dot{\bm{\omega}}
 \;-\; C_s(t)\,\dot{\bm{T}}_s
 \;-\; K_s(t)\,\bm{T}_s,
\label{eq:slosh}
\end{equation}
with each scalar coefficient time-varying as
\begin{equation}
f(t) \;=\; f_{\text{nom}} + f_{\text{amp}}\sin(\omega_f\,t),
\qquad f \in \{A_s,B_s,C_s,K_s\}.
\end{equation}
The constraint \(f_{\text{nom}} > |f_{\text{amp}}|\) keeps \(K_s, C_s\)
positive (negative stiffness or damping would be unphysical). \(\bm{T}_s\)
enters Euler's equation as an additional external torque on the body.

\textbf{No algebraic loop.} \(\dot{\bm{\omega}}\) depends on \(\bm{T}_s\) (a
\emph{state}, not on \(\ddot{\bm{T}}_s\)), and \(\ddot{\bm{T}}_s\) consumes
the just-computed \(\dot{\bm{\omega}}\). The RHS is fully explicit.

\subsection{Coupled Equations of Motion}

The complete state is
\(\bm{x} = [\bm{q};\bm{\omega};\bm{h}_w;\bm{T}_s;\dot{\bm{T}}_s]
        \in \mathbb{R}^{13+n_w}\),
governed by
\begin{align}
\dot{\bm{q}}        &= \tfrac{1}{2}\,\bm{\Omega}(\bm{\omega})\,\bm{q}, \\
\bm{J}\,\dot{\bm{\omega}} &= \bm{u}_{\text{body}}
                            - \bm{\omega}\times(\bm{J}\bm{\omega}+\bm{W}\bm{h}_w), \\
\dot{\bm{h}}_w      &= \bm{\tau}_w, \\
\ddot{\bm{T}}_s     &= -A_s\bm{\omega} - B_s\dot{\bm{\omega}}
                       - C_s\dot{\bm{T}}_s - K_s\bm{T}_s,
\end{align}
with
\(\bm{u}_{\text{body}} = -\bm{W}\bm{\tau}_w + \bm{m}\times\bm{B} + \bm{T}_s\).


% ---------------------------------------------------------------------
\section{Code Architecture Evolution}
% ---------------------------------------------------------------------

The simulation was built incrementally; each step is reproducible from
the conversation log.

\begin{enumerate}\setlength{\itemsep}{4pt}
  \item \textbf{Baseline rigid body + PD controller.}
        A single-file implementation of Euler's equation with a quaternion
        PD law and a hand-written fourth-order Runge--Kutta integrator.
  \item \textbf{Euler-angle initial conditions.}
        Helper \texttt{euler321\_to\_quat()} converts a roll-pitch-yaw
        triple to a scalar-last quaternion (M\&C App.~B.96).
  \item \textbf{Full 3$\times$3 inertia tensor.}
        Replaced the diagonal default with off-diagonal products of inertia
        and added symmetry / positive-definiteness assertions.
  \item \textbf{Reaction-wheel dynamics.}
        Added the four-wheel NASA pyramid configuration matrix \(\bm{W}\),
        pseudoinverse allocation, and the wheel-momentum gyroscopic term in
        Euler's equation.
  \item \textbf{Magnetic momentum unloading.}
        Added a toy LEO B-field, the cross-product unloading law, and per-
        axis dipole saturation. Also introduced PD torque saturation
        \(u_{\max}\).
  \item \textbf{Propellant slosh disturbance.}
        Added the second-order slosh ODE driven by \((\bm{\omega},
        \dot{\bm{\omega}})\), with time-varying scalar coefficients
        (DC~+~sine).
  \item \textbf{Variable rename.}
        Renamed \(\Gamma_s \to T_s\) throughout the code and plot labels
        for consistency with the user's preferred notation.
  \item \textbf{Toggles.}
        Added independent on/off switches for slosh, magnetorquer, and
        reaction wheels. With wheels disabled the PD output drives the body
        directly as an ideal torque actuator.
  \item \textbf{Migration to \texttt{solve\_ivp}.}
        Replaced the hand-rolled RK4 loop with
        \texttt{scipy.integrate.solve\_ivp} (adaptive RK45). Quaternion
        re-normalisation moved from per-step to post-processing.
  \item \textbf{Merge and clean-up.}
        Consolidated the two intermediate files into a single
        \texttt{rigid\_body\_pd.py}. Removed all RK4-specific code paths.
  \item \textbf{Dataclass + presets refactor (v2).}
        Centralised every spacecraft parameter, gain, initial condition,
        time setting, and toggle into a single \texttt{SimConfig} dataclass
        at the top of the file. Added two named preset factories,
        \texttt{mc\_example\_71()} and \texttt{mc\_example\_72()}, matching
        the worked examples in M\&C Sec.~7.
\end{enumerate}


% ---------------------------------------------------------------------
\section{Final Architecture (v2)}
% ---------------------------------------------------------------------

\subsection{The \texttt{SimConfig} Dataclass}

All scenario knobs live in one place:

\begin{lstlisting}
@dataclass
class SimConfig:
    """All satellite parameters, gains, ICs, time, and toggles."""
    # Inertia (M&C Eq. 3.13), kg.m^2
    Ixx: float = 6400.0;  Iyy: float = 4730.0;  Izz: float = 8160.0
    Ixy: float = 76.4;    Ixz: float = 25.6;    Iyz: float = 40.0
    # Wheel geometry
    wheel_tilt_deg: float = 54.7356      # 4-wheel NASA pyramid
    # PD gains and saturation
    Kp: float = 10.0;   Kd: float = 150.0;   u_max: float = 1000.0
    # Magnetorquer
    K_dump: float = 1.0;  m_max: float = 500.0
    # Slosh:  f(t) = nom + amp*sin(omega*t)
    slosh_params: dict = field(default_factory=lambda: {
        'A': {'nom': 0.005, 'amp': 0.0, 'omega': 0.05},
        'B': {'nom': 0.020, 'amp': 0.0, 'omega': 0.05},
        'C': {'nom': 0.023, 'amp': 0.0, 'omega': 0.05},
        'K': {'nom': 0.008, 'amp': 0.0, 'omega': 0.05},
    })
    # Initial conditions (scalar-last)
    q0:    np.ndarray = field(default_factory=lambda:
        (math.sqrt(2)/2) * np.array([1.0, 0.0, 0.0, 1.0]))
    w0:    np.ndarray = field(default_factory=lambda:
        np.array([0.01, 0.01, 0.01]))
    q_des: np.ndarray = field(default_factory=lambda:
        np.array([0.0, 0.0, 0.0, 1.0]))
    Ts0:   np.ndarray = field(default_factory=lambda: np.zeros(3))
    Tsd0:  np.ndarray = field(default_factory=lambda: np.zeros(3))
    # Time
    t_end: float = 20*60.0;  dt_eval: float = 0.05
    # Toggles
    enable_slosh:        bool = False
    enable_magnetorquer: bool = False
    enable_wheels:       bool = True
    # Solver
    method: str = 'RK45';  rtol: float = 1e-8
    atol: float = 1e-10;   max_step: float = 0.5
\end{lstlisting}

\subsection{Named Presets}

\begin{lstlisting}
def mc_example_71() -> SimConfig:
    """Large spacecraft, diagonal J, fast initial tumble."""
    cfg = SimConfig()
    cfg.Ixx, cfg.Iyy, cfg.Izz = 10000.0, 9000.0, 12000.0
    cfg.Ixy = cfg.Ixz = cfg.Iyz = 0.0
    cfg.Kp, cfg.Kd = 50.0, 500.0
    cfg.u_max = 100.0
    q = np.array([0.6853, 0.6953, 0.1531, 0.1531])
    cfg.q0 = q / np.linalg.norm(q)
    cfg.w0 = np.array([0.53, 0.53, 0.053])
    cfg.t_end = 300.0
    return cfg

def mc_example_72() -> SimConfig:
    """Smaller spacecraft with products of inertia; gentle IC."""
    return SimConfig()    # SimConfig() defaults already match Example 7.2
\end{lstlisting}

\subsection{Usage Pattern}

\begin{lstlisting}
cfg = mc_example_72()
cfg.enable_slosh = True            # in-place tweak
cfg.Kp, cfg.Kd = 15.0, 200.0
results = simulate(cfg)            # one entry point
plot_results(*results)
plot_euler_and_rates(results[0], results[1], results[3])
plt.show()
\end{lstlisting}


% ---------------------------------------------------------------------
\section{Validation: Automated Smoke Test}
% ---------------------------------------------------------------------

A separate file \texttt{test\_rigid\_body\_pd\_v2.py} exercises the
simulation under every toggle combination using the non-interactive
\texttt{Agg} matplotlib backend. All eight test cases pass.

\begin{center}
\begin{tabular}{lc}
\toprule
\textbf{Test} & \textbf{Result} \\
\midrule
Euler $\leftrightarrow$ quaternion roundtrip (4 angle triples) & Pass \\
Example 7.1 preset, wheels only                                 & Pass \\
Example 7.2 preset, wheels only                                 & Pass \\
Example 7.2 full stack (slosh + mag + wheels)                   & Pass \\
Example 7.2 with wheels OFF (ideal actuator)                    & Pass \\
Example 7.2 with all subsystems OFF                             & Pass \\
Example 7.2 with overridden gains (in-place mutation)           & Pass \\
Plot functions (both figures)                                   & Pass \\
\bottomrule
\end{tabular}
\end{center}

\noindent Per-trajectory invariants checked at every output sample:
\begin{itemize}\setlength{\itemsep}{1pt}
  \item Output tuple length = 11; all array shapes match the expected
        \((N,4)\), \((N,3)\), \((N,n_w)\).
  \item No NaNs in any state trajectory.
  \item Unit-norm quaternion preserved to machine precision
        (\(\max\bigl|\,\|\bm{q}\|-1\,\bigr| = 2.22\times10^{-16}\)).
  \item \(|\bm{\omega}|\) monotonically non-increasing in every settling case.
  \item When a toggle is OFF, the corresponding telemetry stays at
        exactly zero (\(\bm{h}_w\), \(\bm{\tau}_w\), \(\bm{m}\),
        \(\bm{\tau}_{\text{mag}}\), \(\bm{T}_s\)).
\end{itemize}


% ---------------------------------------------------------------------
\section{Companion LaTeX Physics Summary}
% ---------------------------------------------------------------------

A separate file \texttt{attitude\_dynamics\_model.tex} was produced for
direct import into the project's Overleaf write-up. It contains a
self-contained section with every equation above plus extended
discussion, fully cross-referenced via \texttt{\textbackslash label}.
The \texttt{\textbackslash cite\{markley2014\}} entry should resolve to
Markley and Crassidis using the BibTeX stanza included at the top of that
file.


% ---------------------------------------------------------------------
\section{Proposed Extension: NARX Slosh Surrogate}
% ---------------------------------------------------------------------

A design (no code) was discussed for replacing the analytical slosh ODE
\eqref{eq:slosh} with a learned Nonlinear Autoregressive Network with
eXogenous Inputs (NARX). The map to learn is
\begin{equation}
\hat{T}_s(k) \;=\; f_\theta\bigl(
  T_s(k{-}1),\dots,T_s(k{-}n_y),\;
  \omega(k),\dots,\omega(k{-}n_u),\;
  \dot\omega(k),\dots,\dot\omega(k{-}n_u)
\bigr),
\end{equation}
optionally extended to a \emph{parametric} surrogate by also feeding the
four slosh coefficients \((A_s, B_s, C_s, K_s)\). The plan, in order:

\begin{enumerate}\setlength{\itemsep}{2pt}
  \item Add a helper \texttt{reconstruct\_w\_dot()} that recomputes
        \(\dot{\bm{\omega}}\) exactly from the saved \texttt{simulate()}
        outputs (no finite differences).
  \item Implement \texttt{sample\_random\_config(rng)} that draws
        \texttt{SimConfig} instances with randomised initial conditions,
        gains, inertia perturbations, and (optionally) slosh coefficients.
  \item Generate 100--500 trajectories, each $\sim$600\,s, downsampled to
        \(\Delta t \approx 0.5\,\mathrm{s}\) and stacked across all three
        body axes (the slosh dynamics are axis-decoupled, so one scalar
        NARX trained on stacked per-axis data is the right architecture).
  \item Build sliding-window features with lag orders
        \(n_y \approx n_u \approx 10\)--\(20\) samples.
  \item Train a small MLP (e.g. 64-64 \texttt{tanh}) with
        \texttt{sklearn.neural\_network.MLPRegressor} as a baseline, then
        promote to PyTorch if needed.
  \item Validate in two modes:
        \emph{one-step prediction} (teacher-forced lags) for training-loss
        sanity, and \emph{free-running rollout} (model feeds its own past
        predictions) for honest dynamic accuracy.
  \item Split by \emph{whole trajectory} (70/15/15) to avoid leakage from
        adjacent correlated samples within one simulation.
  \item Optional integration: add \texttt{cfg.use\_narx\_slosh = True}
        and \texttt{cfg.narx\_model = trained\_model} to the dataclass,
        wire it into \texttt{dynamics\_rhs}, and compare closed-loop
        traces against the analytical slosh ground truth on held-out
        initial conditions.
\end{enumerate}


% ---------------------------------------------------------------------
\section{Deliverables}
% ---------------------------------------------------------------------

\begin{center}
\begin{tabular}{ll}
\toprule
\textbf{File} & \textbf{Purpose} \\
\midrule
\texttt{rigid\_body\_pd.py}            & Merged simulation, \texttt{solve\_ivp}-based \\
\texttt{rigid\_body\_pd\_v2.py}        & Dataclass + presets refactor (recommended) \\
\texttt{test\_rigid\_body\_pd\_v2.py}  & Smoke test (8 cases, all pass) \\
\texttt{attitude\_dynamics\_model.tex} & Physics section for Overleaf \\
\texttt{session\_summary.tex}          & This document \\
\bottomrule
\end{tabular}
\end{center}


% ---------------------------------------------------------------------
\section{Conventions Used Throughout}
% ---------------------------------------------------------------------

\begin{itemize}\setlength{\itemsep}{1pt}
  \item \textbf{Quaternion:} scalar-last, \(\bm{q}=[q_1, q_2, q_3, q_4]\)
        with \(q_4\) the scalar part.
  \item \textbf{Euler sequence:} 3-2-1 (Z-Y-X, yaw-pitch-roll) for human
        I/O; quaternions everywhere internally.
  \item \textbf{Inertia sign:} off-diagonal entries are stored as positive
        products of inertia; the negative signs live in the \(\bm{J}\)
        builder (M\&C Eq.~3.13).
  \item \textbf{Frames:} body frame throughout; \(\bm{B}\) and
        \(\bm{h}_w^{\,b}\) are body-frame projections.
  \item \textbf{Slosh torque symbol:} \(\bm{T}_s\) (code: \texttt{Ts}),
        with rate \(\dot{\bm{T}}_s\) (\texttt{Ts\_dot}). Renamed from the
        original \(\Gamma_s\).
\end{itemize}


% ---------------------------------------------------------------------
\begin{thebibliography}{1}
% ---------------------------------------------------------------------

\bibitem{markley2014}
F.~L.~Markley and J.~L.~Crassidis,
\emph{Fundamentals of Spacecraft Attitude Determination and Control},
Space Technology Library, Vol.~33,
Springer, New York, NY, 2014.
\href{https://doi.org/10.1007/978-1-4939-0802-8}{doi:10.1007/978-1-4939-0802-8}.

\end{thebibliography}

\end{document}
